{ %description
    \small
	\selectlanguage{russian}
	\raggedright
	\thispagestyle{empty} % delete number of  page
	\par Соснин В.В., Балакшин П.В., Шилко Д.С. Введение в параллельные вычисления. -- СПб: Университет ИТМО, 2020. Электронные исходники пособия.
	\par\textit{Редакционное примечание от 24.10.2026.} Исходное авторство и год сохранены. Объём нового PDF не приравнивается к непроверенной пометке «51 с.» из архива. Добавлены датированные уточнения и глава Concurrency Theatre; это не заявление о новом официальном издании Университета ИТМО.
	\newline
	\newline
	\par В пособии излагаются основные понятия и определения теории параллельных вычислений. Рассматриваются основные принципы построения программ на языке «Си» для многоядерных и многопроцессорных вычислительных комплексов с общей памятью. Предлагается набор заданий для проведения лабораторных и практических занятий.
	\newline
	\newline
	\par Учебное пособие предназначено для студентов, обучающихся по магистерским программам направления «09.04.04 – Программная инженерия», «09.04.01 –Информатика и вычислительная техника», и может быть использовано выпускниками (бакалаврами и магистрантами) при написании выпускных квалификационных работ, связанных с проектированием и исследованием многоядерных и многопроцессорных вычислительных комплексов.
	\newline
	\newline
	% An unverified archival approval statement is not reproduced as a current approval.
	\vspace*{\fill} %text in the bottom of page
	\begin{flushright}
		\includegraphics[width=9cm, height=1.5cm]{ITMOLogo2}
	\end{flushright}
	\par\textit{Следующие сведения об университете воспроизведены из описания 2020 года.}
	\par Университет ИТМО – ведущий вуз России в области информационных и фотонных технологий, один из немногих российских вузов, получивших в 2009 году статус национального исследовательского университета. С 2013 года Университет ИТМО – участник программы повышения конкурентоспособности российских университетов среди ведущих мировых научно-образовательных центров, известной как проект «5 в 100». Цель Университета ИТМО – становление исследовательского университета мирового уровня, предпринимательского по типу, ориентированного на интернационализацию всех направлений деятельности.
	\begin{flushright}
		\copyright\space Университет ИТМО, 2020
		\par\copyright\space Соснин В.В., Балакшин П.В., Шилко Д.С., 2020
	\end{flushright}
}